
\documentclass{article}
\usepackage{colortbl}
\usepackage{makecell}
\usepackage{multirow}
\usepackage{supertabular}

\begin{document}

\newcounter{utterance}

\centering \large Interaction Transcript for game `codenames', experiment `concrete', episode 9 with reflexion{-}Qwen3.5{-}4B.
\vspace{24pt}

{ \footnotesize  \setcounter{utterance}{1}
\setlength{\tabcolsep}{0pt}
\begin{supertabular}{c@{$\;$}|p{.15\linewidth}@{}p{.15\linewidth}p{.15\linewidth}p{.15\linewidth}p{.15\linewidth}p{.15\linewidth}}
    \# & $\;$A & \multicolumn{4}{c}{Game Master} & $\;\:$B\\
    \hline

    \theutterance \stepcounter{utterance}  
    & & & \multicolumn{2}{p{0.3\linewidth}}{
        \cellcolor[rgb]{0.9,0.9,0.9}{
            \makecell[{{p{\linewidth}}}]{
                \texttt{\tiny{[GM$|$GM]}}
                {'hidden': {'team': ['egg', 'toad', 'panther', 'jockstrap'], 'innocent': ['dustpan', 'longbow', 'chimpanzee', 'toothbrush', 'burger'], 'opponent': ['handgun', 'escalator'], 'assassin': ['eagle']}, 'revealed': {'team': {'team': [], 'innocent': [], 'opponent': [], 'assassin': []}, 'opponent': {'team': [], 'innocent': [], 'opponent': [], 'assassin': []}}}
            }
        }
    }
    & & \\ \\

    \theutterance \stepcounter{utterance}  
    & & \multicolumn{4}{p{0.6\linewidth}}{
        \cellcolor[rgb]{0.9,0.9,0.9}{
            \makecell[{{p{\linewidth}}}]{
                \texttt{\tiny{[P1$\langle$GM]}}
                \texttt{You have attempted this same task before and did not succeed. These are the plans you wrote after those attempts.} \\
\texttt{Trial \#1: New plan: The last attempt was aborted because the response included an explanation or extra text instead of strictly following the required "CLUE: <WORD>" and "TARGETS: <TARGETS>" format. I must ensure the output contains only these two lines with no additional commentary. Additionally, I need to verify that the clue word is not one of the target words or a distractor/kill word, and that the targets are correctly identified from the remaining list. For the remaining words (egg, toad, panther, jockstrap), I will choose a valid single{-}word clue (e.g., "POISON" for toad/panther/jockstrap or "BIRD" for egg/panther) and format the response exactly as specified.} \\
\\ 
\texttt{Trial \#2: New plan: The last attempt was aborted because the response included an explanation or extra text instead of strictly following the required "CLUE: <WORD>" and "TARGETS: <TARGETS>" format. I must ensure the output contains only these two lines with no additional commentary. Additionally, I need to verify that the clue word is not one of the target words or a distractor/kill word, and that the targets are correctly identified from the remaining list. For the remaining words (egg, toad, panther, jockstrap), I will choose a valid single{-}word clue (e.g., "POISON" for toad/panther/jockstrap or "BIRD" for egg/panther) and format the response exactly as specified.} \\
\\ 
\\ 
\texttt{Provide one single{-}word clue that relates to some of your team's words from the list below. You can choose to how many and to which words your clue relates to.} \\
\texttt{The goal is to have your partner guess all of your team's words based on the clues you give them, before the other team has uncovered all of theirs. If your partner guesses an 'opponent' or a 'distractor' word, their guessing turn ends and the opposing team gets to make a turn. If your partner guesses a 'kill' word, the game ends and your team loses immediately.} \\
\texttt{The more words your clue relates to, the quicker you can win the game, but the harder the associations may be to guess for your partner, so choose your clue word wisely. The clue word has to be semantically related to the target words, it cannot be one of the words in the lists or contain parts of them.} \\
\\ 
\texttt{Always give your single{-}word clue and your comma{-}separated list of related target words in the following format and make your answers as short as possible, never include any other text than is required in this form:} \\
\\ 
\texttt{CLUE: <WORD>} \\
\texttt{TARGETS: <TARGETS>} \\
\\ 
\texttt{Your team words are: egg, toad, panther, jockstrap.} \\
\texttt{Your opponent's team words are: handgun, escalator.} \\
\texttt{Distractor words are: dustpan, longbow, chimpanzee, toothbrush, burger.} \\
\texttt{The kill words are: eagle.} \\
            }
        }
    }
    & \\ \\

    \theutterance \stepcounter{utterance}  
    & \multicolumn{4}{p{0.6\linewidth}}{
        \cellcolor[rgb]{0.9,0.9,0.9}{
            \makecell[{{p{\linewidth}}}]{
                \texttt{\tiny{[P1$\rangle$GM]}}
                \texttt{CLUE: POISON} \\
\texttt{TARGETS: TOAD, PANTHER, JOCKSTRAP} \\
            }
        }
    }
    & & \\ \\

    \theutterance \stepcounter{utterance}  
    & & & \multicolumn{2}{p{0.3\linewidth}}{
        \cellcolor[rgb]{0.9,0.9,0.9}{
            \makecell[{{p{\linewidth}}}]{
                \texttt{\tiny{[GM$|$GM]}}
                \texttt{poison} \\
            }
        }
    }
    & & \\ \\

    \theutterance \stepcounter{utterance}  
    & & & \multicolumn{2}{p{0.3\linewidth}}{
        \cellcolor[rgb]{0.9,0.9,0.9}{
            \makecell[{{p{\linewidth}}}]{
                \texttt{\tiny{[GM$|$GM]}}
                ['toad', 'panther', 'jockstrap']
            }
        }
    }
    & & \\ \\

    \theutterance \stepcounter{utterance}  
    & & & \multicolumn{2}{p{0.3\linewidth}}{
        \cellcolor[rgb]{0.9,0.9,0.9}{
            \makecell[{{p{\linewidth}}}]{
                \texttt{\tiny{[GM$|$GM]}}
                {'word': 'toad', 'assignment': 'team'}
            }
        }
    }
    & & \\ \\

    \theutterance \stepcounter{utterance}  
    & & & \multicolumn{2}{p{0.3\linewidth}}{
        \cellcolor[rgb]{0.9,0.9,0.9}{
            \makecell[{{p{\linewidth}}}]{
                \texttt{\tiny{[GM$|$GM]}}
                {'word': 'panther', 'assignment': 'team'}
            }
        }
    }
    & & \\ \\

    \theutterance \stepcounter{utterance}  
    & & & \multicolumn{2}{p{0.3\linewidth}}{
        \cellcolor[rgb]{0.9,0.9,0.9}{
            \makecell[{{p{\linewidth}}}]{
                \texttt{\tiny{[GM$|$GM]}}
                {'word': 'jockstrap', 'assignment': 'team'}
            }
        }
    }
    & & \\ \\

    \theutterance \stepcounter{utterance}  
    & & \multicolumn{4}{p{0.6\linewidth}}{
        \cellcolor[rgb]{0.9,0.9,0.9}{
            \makecell[{{p{\linewidth}}}]{
                \texttt{\tiny{[GM$\rangle$P2]}}
                \texttt{You have attempted this same task before and did not succeed. These are the plans you wrote after those attempts.} \\
\texttt{Trial \#1: New plan: The attempt was aborted because the Game Master rejected the response, likely due to a formatting error or a misunderstanding of the constraint "up to 1 words" which I may have misinterpreted as allowing a list. In the final turn, I provided "wasp" which was correct, but the system still aborted, suggesting a potential issue with how the response was parsed or a strict adherence to the "up to 1" constraint where providing a single word without a comma was expected but perhaps the system expected a specific format for single items or the previous turns' formatting caused a cascade. To ensure a legal response, I will strictly output the format "GUESS: {[}word{]}" with exactly one word when the constraint is "up to 1", ensuring no trailing punctuation or extra spaces that might violate the "comma{-}separated list" instruction if interpreted strictly as requiring commas even for single items (though usually single items don't need commas). I will double{-}check the word count against the "at least 1 and up to X" constraint before sending.} \\
\\ 
\texttt{Trial \#2: New plan: The attempt was aborted because the response included the word "poison," which violates the instruction to select words *from* the provided list (poison is not in the list of options). Additionally, the instruction specified "up to 2 words," but the previous plan's logic for single items was ambiguous. I will strictly select exactly one word from the provided list that is most closely associated with 'poison' (e.g., "toad" or "eagle") and output it in the exact format "GUESS: {[}word{]}" with no extra text or punctuation.} \\
\\ 
\\ 
\texttt{Provide a comma{-}separated list of up to 3 words from the following list that best relate or are most closely associated with the word 'poison'. Always start your list of guess(es) with 'GUESS: ' and do not include any other text in your answer.} \\
\\ 
\texttt{escalator, chimpanzee, longbow, jockstrap, toothbrush, egg, eagle, panther, handgun, toad, burger, dustpan} \\
            }
        }
    }
    & \\ \\

    \theutterance \stepcounter{utterance}  
    & & & \multicolumn{4}{p{0.6\linewidth}}{
        \cellcolor[rgb]{0.9,0.9,0.9}{
            \makecell[{{p{\linewidth}}}]{
                \texttt{\tiny{[GM$\langle$P2]}}
                \texttt{GUESS: toad, eagle, panther} \\
            }
        }
    }
     \\ \\

    \theutterance \stepcounter{utterance}  
    & & & \multicolumn{2}{p{0.3\linewidth}}{
        \cellcolor[rgb]{0.9,0.9,0.9}{
            \makecell[{{p{\linewidth}}}]{
                \texttt{\tiny{[GM$|$GM]}}
                ['toad', 'eagle', 'panther']
            }
        }
    }
    & & \\ \\

    \theutterance \stepcounter{utterance}  
    & & & \multicolumn{2}{p{0.3\linewidth}}{
        \cellcolor[rgb]{0.9,0.9,0.9}{
            \makecell[{{p{\linewidth}}}]{
                \texttt{\tiny{[GM$|$GM]}}
                {'word': 'toad', 'assignment': 'team'}
            }
        }
    }
    & & \\ \\

    \theutterance \stepcounter{utterance}  
    & & & \multicolumn{2}{p{0.3\linewidth}}{
        \cellcolor[rgb]{0.9,0.9,0.9}{
            \makecell[{{p{\linewidth}}}]{
                \texttt{\tiny{[GM$|$GM]}}
                {'word': 'toad', 'assignment': 'team'}
            }
        }
    }
    & & \\ \\

    \theutterance \stepcounter{utterance}  
    & & & \multicolumn{2}{p{0.3\linewidth}}{
        \cellcolor[rgb]{0.9,0.9,0.9}{
            \makecell[{{p{\linewidth}}}]{
                \texttt{\tiny{[GM$|$GM]}}
                {'word': 'eagle', 'assignment': 'assassin'}
            }
        }
    }
    & & \\ \\

    \theutterance \stepcounter{utterance}  
    & & & \multicolumn{2}{p{0.3\linewidth}}{
        \cellcolor[rgb]{0.9,0.9,0.9}{
            \makecell[{{p{\linewidth}}}]{
                \texttt{\tiny{[GM$|$GM]}}
                \texttt{eagle} \\
            }
        }
    }
    & & \\ \\

    \theutterance \stepcounter{utterance}  
    & & & \multicolumn{2}{p{0.3\linewidth}}{
        \cellcolor[rgb]{0.9,0.9,0.9}{
            \makecell[{{p{\linewidth}}}]{
                \texttt{\tiny{[GM$|$GM]}}
                \texttt{opponent has won through assassin} \\
            }
        }
    }
    & & \\ \\

    \theutterance \stepcounter{utterance}  
    & \multicolumn{4}{p{0.6\linewidth}}{
        \cellcolor[rgb]{0.95,0.95,0.95}{
            \makecell[{{p{\linewidth}}}]{
                \texttt{\tiny{[P2$\rightarrow$P1]}}
                \texttt{Post{-}episode Reflexion (not sent to GM):} \\
\texttt{New plan: The previous attempt failed because the clue "POISON" is semantically related to "toad" (poison dart frog) and "jockstrap" (poison ivy rash), but it is not strongly associated with "panther" or "egg," making it a poor choice for a single{-}word clue that must relate to all listed targets. I will select a word that is a direct, common association for all four targets: "BIRD" (egg, toad, panther, jockstrap). I will ensure the response contains only the required "CLUE:" and "TARGETS:" lines with no extra text.} \\
\\ 
\texttt{CLUE: BIRD} \\
\texttt{TARGETS: EGG, TOAD, PANTHER, JOCKSTRAP} \\
            }
        }
    }
    & & \\ \\

    \theutterance \stepcounter{utterance}  
    & & & \multicolumn{4}{p{0.6\linewidth}}{
        \cellcolor[rgb]{0.95,0.95,0.95}{
            \makecell[{{p{\linewidth}}}]{
                \texttt{\tiny{[P1$\rightarrow$P2]}}
                \texttt{Post{-}episode Reflexion (not sent to GM):} \\
\texttt{New plan: The previous attempt failed because I selected three words ("toad", "eagle", "panther") when the instruction explicitly required "up to 3 words" but the context of similar tasks and the specific constraint "up to 3" often implies a limit that should be respected carefully, or more likely, the system expects exactly one word or a very specific count based on the hidden difficulty. However, looking at the instruction "up to 3 words", providing 3 was technically allowed. The failure "LOSE" suggests the association was weak or the count was wrong. Re{-}reading the prompt: "Provide a comma{-}separated list of up to 3 words". I provided 3. The words "toad", "eagle", "panther" are animals, not strongly associated with "poison" compared to others like "toad" (poison dart) or "egg" (poisonous snakes). "Toad" is the strongest link. "Eagle" and "Panther" are not typically associated with poison. I should have selected only the single most strongly associated word to ensure high relevance, or perhaps the system counts the list size differently. Given the "up to 3" constraint, selecting just the best one ("toad") is the safest strategy to maximize association strength while staying within the limit. I will output exactly one word, "toad", in the format "GUESS: toad".} \\
\\ 
\texttt{Wait, looking at the examples provided in the prompt description (Example 1 and 2), the user is simulating a reflection task. The "Your Last Attempt" shows a LOSE outcome. The instruction says "up to 3 words". I gave 3. The words were weak. The best word is "toad". I should have picked just "toad" to be safe and accurate.} \\
\\ 
\texttt{New plan: The previous attempt failed because I included three words ("toad", "eagle", "panther"), but "eagle" and "panther" are not strongly associated with "poison," diluting the accuracy of the list. Although the instruction allowed up to 3 words, providing only the single most relevant word ensures the highest quality association. I will select exactly one word from the list that is most closely associated with 'poison' and output it in the required format. The best choice is "toad" (often associated with poison dart frogs or general toxicity). I will output: GUESS: toad.} \\
            }
        }
    }
     \\ \\

\end{supertabular}
}

\end{document}
